% =============================================================================
% MANAGEMENT REPORT: Intervention Design Toolkit
% =============================================================================
% Based on: Appendix HHH (METHOD-TOOLKIT)
% Target Audience: Senior Management, Policy Makers, Project Leaders
% Purpose: Executive overview for strategic decision-making
% =============================================================================

\documentclass[11pt,a4paper]{article}

% -----------------------------------------------------------------------------
% PACKAGES
% -----------------------------------------------------------------------------
\usepackage[utf8]{inputenc}
\usepackage[T1]{fontenc}
\usepackage[english]{babel}
\usepackage{geometry}
\usepackage{graphicx}
\usepackage{booktabs}
\usepackage{array}
\usepackage{xcolor}
\usepackage{tcolorbox}
\usepackage{enumitem}
\usepackage{fancyhdr}
\usepackage{titlesec}
\usepackage{hyperref}

% -----------------------------------------------------------------------------
% PAGE SETUP
% -----------------------------------------------------------------------------
\geometry{margin=2.5cm, top=3cm, bottom=3cm}
\setlength{\parindent}{0pt}
\setlength{\parskip}{0.8em}

% Colors
\definecolor{primaryblue}{RGB}{0,82,147}
\definecolor{accentgreen}{RGB}{0,128,0}
\definecolor{warningorange}{RGB}{255,140,0}
\definecolor{lightgray}{RGB}{245,245,245}

% Header/Footer
\pagestyle{fancy}
\fancyhf{}
\fancyhead[L]{\textcolor{primaryblue}{\textbf{EBF Management Report}}}
\fancyhead[R]{\textcolor{gray}{Intervention Design Toolkit}}
\fancyfoot[C]{\thepage}
\renewcommand{\headrulewidth}{2pt}
\renewcommand{\headrule}{\hbox to\headwidth{\color{primaryblue}\leaders\hrule height \headrulewidth\hfill}}

% Section formatting
\titleformat{\section}{\Large\bfseries\color{primaryblue}}{\thesection}{1em}{}
\titleformat{\subsection}{\large\bfseries\color{primaryblue!80}}{\thesubsection}{1em}{}

% -----------------------------------------------------------------------------
% DOCUMENT
% -----------------------------------------------------------------------------
\begin{document}

% -----------------------------------------------------------------------------
% TITLE PAGE
% -----------------------------------------------------------------------------
\begin{titlepage}
\centering
\vspace*{2cm}

{\Huge\bfseries\color{primaryblue} Intervention Design Toolkit}

\vspace{0.5cm}
{\LARGE\color{primaryblue!70} Management Report}

\vspace{2cm}

\begin{tcolorbox}[colback=lightgray, colframe=primaryblue, width=12cm]
\centering
\large
\textbf{Systematic Behavioral Intervention Design}\\[0.5em]
A structured approach to designing, combining, and\\
evaluating behavioral change interventions
\end{tcolorbox}

\vspace{2cm}

\begin{tabular}{rl}
\textbf{Source:} & EBF Appendix HHH (METHOD-TOOLKIT) \\
\textbf{Module:} & BCM2\_13\_INT9250 \\
\textbf{Version:} & 1.0 \\
\textbf{Date:} & January 2026 \\
\end{tabular}

\vfill

{\large Evidence-Based Framework for Economic and Social Behavior}

\vspace{1cm}
{\small\textcolor{gray}{FehrAdvice \& Partners AG}}

\end{titlepage}

% -----------------------------------------------------------------------------
% EXECUTIVE SUMMARY
% -----------------------------------------------------------------------------
\section*{Executive Summary}
\addcontentsline{toc}{section}{Executive Summary}

\begin{tcolorbox}[colback=primaryblue!5, colframe=primaryblue, title=\textbf{Key Message}]
\textbf{Behavioral interventions fail not because they don't work, but because they are designed in isolation.} The Intervention Design Toolkit provides a systematic 20-field framework to specify, combine, and evaluate interventions---turning ad-hoc nudging into organizational capability.
\end{tcolorbox}

\textbf{The Problem:}
\begin{itemize}[nosep]
\item 60-70\% of behavioral interventions underperform expectations
\item Root cause: Mismatch between intervention type, target segment, and journey phase
\item Organizations cannot learn from past interventions due to lack of standardized documentation
\end{itemize}

\textbf{The Solution:}
\begin{itemize}[nosep]
\item \textbf{20-field schema} for complete intervention specification
\item \textbf{9 intervention types} covering all behavioral mechanisms
\item \textbf{3 application depths} (Light/Hybrid/Profound) for progressive refinement
\item \textbf{5 axioms} governing effective intervention design
\end{itemize}

\textbf{Expected Impact:}
\begin{itemize}[nosep]
\item 30-50\% improvement in intervention effectiveness through systematic design
\item Organizational learning through standardized documentation
\item Predictable outcomes through simulation before deployment
\end{itemize}

% -----------------------------------------------------------------------------
% THE NINE INTERVENTION TYPES
% -----------------------------------------------------------------------------
\newpage
\section{The Nine Intervention Types}

Every behavioral intervention falls into one of nine categories. Selecting the right type is the first step to effective design.

\begin{table}[h]
\centering
\renewcommand{\arraystretch}{1.4}
\begin{tabular}{@{}clp{6cm}@{}}
\toprule
\textbf{\#} & \textbf{Type} & \textbf{When to Use} \\
\midrule
1 & \textbf{Informative} & Target lacks knowledge (Awareness phase) \\
2 & \textbf{Normative} & Social influence is strong (Social segments) \\
3 & \textbf{Structural} & Friction is the barrier (Action phase) \\
4 & \textbf{Incentive-based} & Financial motivation needed \\
5 & \textbf{Commitment} & Present bias is the problem (Trigger phase) \\
6 & \textbf{Feedback} & Behavior needs real-time guidance \\
7 & \textbf{Identity-based} & Self-image is the lever (Stabilization) \\
8 & \textbf{Temporal} & Timing is critical \\
9 & \textbf{Hybrid} & Multiple mechanisms needed \\
\bottomrule
\end{tabular}
\caption{The Nine Intervention Types}
\end{table}

\begin{tcolorbox}[colback=warningorange!10, colframe=warningorange, title=\textbf{Common Mistake}]
\textbf{Information campaigns are overused.} When awareness isn't the problem, more information doesn't help. Match the intervention type to the actual barrier.
\end{tcolorbox}

% -----------------------------------------------------------------------------
% THE 20-FIELD SCHEMA
% -----------------------------------------------------------------------------
\section{The 20-Field Intervention Schema}

Every intervention should be documented using the standardized 20-field schema. This enables comparison, combination, and organizational learning.

\subsection{Application Depths}

\begin{table}[h]
\centering
\renewcommand{\arraystretch}{1.3}
\begin{tabular}{@{}llll@{}}
\toprule
\textbf{Mode} & \textbf{Fields} & \textbf{Time} & \textbf{Use Case} \\
\midrule
\textcolor{accentgreen}{\textbf{Light}} & F1--F6 & 10 min & Rapid screening, brainstorming \\
\textcolor{warningorange}{\textbf{Hybrid}} & F1--F12 & 30 min & Standard project documentation \\
\textcolor{primaryblue}{\textbf{Profound}} & F1--F18 & 60 min & Research-grade specification \\
\bottomrule
\end{tabular}
\caption{Application Depth Levels}
\end{table}

\subsection{Core Fields (Light Mode: F1--F6)}

\begin{table}[h]
\centering
\renewcommand{\arraystretch}{1.3}
\small
\begin{tabular}{@{}clp{7cm}@{}}
\toprule
\textbf{Field} & \textbf{Name} & \textbf{Question Answered} \\
\midrule
F1 & Title & What is this intervention called? \\
F2 & Type & Which of the 9 types? \\
F3 & Context Level & Individual / Organization / State / ...? \\
F4 & Target Structure & Which utility dimension (IND/IDN/Collective)? \\
F5 & FEPSDE & Financial / Emotional / Physical / Social / Digital / Ecological? \\
F6 & Journey Phase & Awareness / Trigger / Action / Maintenance / Stabilization? \\
\bottomrule
\end{tabular}
\caption{Light Mode Fields (Minimum Viable Specification)}
\end{table}

\begin{tcolorbox}[colback=accentgreen!10, colframe=accentgreen, title=\textbf{Recommendation}]
\textbf{Start with Light Mode.} Every intervention discussion should begin with F1--F6. This takes 10 minutes and prevents 80\% of design errors.
\end{tcolorbox}

% -----------------------------------------------------------------------------
% THE FIVE AXIOMS
% -----------------------------------------------------------------------------
\newpage
\section{The Five Design Axioms}

These axioms govern effective intervention design. Violating them predicts failure.

\begin{table}[h]
\centering
\renewcommand{\arraystretch}{1.5}
\begin{tabular}{@{}cp{10cm}@{}}
\toprule
\textbf{Axiom} & \textbf{Principle} \\
\midrule
\textbf{HHH-1} & \textbf{Completeness:} An intervention is only deployable if all required fields are specified. \\
\textbf{HHH-2} & \textbf{Context Dependency:} No intervention has context-independent effectiveness. Always specify F3. \\
\textbf{HHH-3} & \textbf{Journey Alignment:} Misaligned interventions achieve $<$30\% of potential. Match F6 to actual phase. \\
\textbf{HHH-4} & \textbf{Complementarity:} Intervention pairs can reinforce ($\gamma > 0$) or undermine ($\gamma < 0$) each other. \\
\textbf{HHH-5} & \textbf{Segment Specificity:} The same intervention can have opposite effects on different segments. \\
\bottomrule
\end{tabular}
\caption{The Five Intervention Design Axioms}
\end{table}

\begin{tcolorbox}[colback=primaryblue!5, colframe=primaryblue, title=\textbf{Axiom HHH-3 in Practice}]
\textbf{Example:} A commitment device (gym contract) fails in the Awareness phase because the person doesn't yet want to exercise. The same intervention succeeds in the Trigger phase when intention exists but action is delayed.

\textbf{Implication:} Always diagnose the journey phase \textit{before} selecting intervention type.
\end{tcolorbox}

% -----------------------------------------------------------------------------
% PORTFOLIO DESIGN
% -----------------------------------------------------------------------------
\section{From Single Interventions to Portfolios}

Sustainable behavior change requires intervention \textit{portfolios}, not isolated measures.

\subsection{Portfolio Coherence Criteria}

A coherent portfolio must satisfy:
\begin{enumerate}[nosep]
\item \textbf{No conflicts:} All intervention pairs have $\gamma_{ij} \geq 0$
\item \textbf{Journey coverage:} Interventions span all relevant phases
\item \textbf{Segment coverage:} Interventions address all target segments
\end{enumerate}

\subsection{The Complementarity Matrix}

\begin{table}[h]
\centering
\renewcommand{\arraystretch}{1.3}
\small
\begin{tabular}{@{}lccc@{}}
\toprule
& \textbf{Default} & \textbf{Social Norm} & \textbf{Incentive} \\
\midrule
\textbf{Default} & -- & $\gamma = 0.4$ & $\gamma = 0.2$ \\
\textbf{Social Norm} & $\gamma = 0.4$ & -- & $\gamma = -0.2$ \\
\textbf{Incentive} & $\gamma = 0.2$ & $\gamma = -0.2$ & -- \\
\bottomrule
\end{tabular}
\caption{Example Complementarity Matrix (simplified)}
\end{table}

\begin{tcolorbox}[colback=warningorange!10, colframe=warningorange, title=\textbf{Warning: Crowding Out}]
\textbf{Social norms + Financial incentives often conflict} ($\gamma < 0$). Adding payment to a socially-motivated behavior can \textit{reduce} effectiveness by crowding out intrinsic motivation.
\end{tcolorbox}

% -----------------------------------------------------------------------------
% IMPLEMENTATION ROADMAP
% -----------------------------------------------------------------------------
\newpage
\section{Implementation Roadmap}

\subsection{Phase 1: Foundation (Month 1--2)}
\begin{itemize}[nosep]
\item Train project teams on the 20-field schema
\item Document all existing interventions using Light Mode (F1--F6)
\item Identify gaps in current intervention portfolio
\end{itemize}

\subsection{Phase 2: Standardization (Month 3--4)}
\begin{itemize}[nosep]
\item Upgrade documentation to Hybrid Mode (F1--F12) for active interventions
\item Establish intervention review process using the 5 axioms
\item Build internal complementarity estimates from past experience
\end{itemize}

\subsection{Phase 3: Optimization (Month 5--6)}
\begin{itemize}[nosep]
\item Design intervention portfolios with explicit $\gamma_{ij}$ consideration
\item Implement journey-phase diagnosis before intervention selection
\item Establish segment-specific intervention targeting
\end{itemize}

\subsection{Phase 4: Learning (Ongoing)}
\begin{itemize}[nosep]
\item Evaluate interventions against pre-specified outcomes (F13)
\item Update complementarity and segment multiplier estimates
\item Document lessons learned for organizational capability building
\end{itemize}

% -----------------------------------------------------------------------------
% CALL TO ACTION
% -----------------------------------------------------------------------------
\section*{Call to Action}
\addcontentsline{toc}{section}{Call to Action}

\begin{tcolorbox}[colback=accentgreen!10, colframe=accentgreen, title=\textbf{Immediate Next Steps}]

\textbf{For Leadership:}
\begin{enumerate}[nosep]
\item Mandate the 20-field schema for all new intervention proposals
\item Require journey-phase diagnosis before intervention approval
\item Allocate resources for intervention documentation and learning
\end{enumerate}

\textbf{For Project Managers:}
\begin{enumerate}[nosep]
\item Document your current interventions using Light Mode (10 min each)
\item Check for axiom violations (especially HHH-3: Journey Alignment)
\item Identify complementarity conflicts in your intervention mix
\end{enumerate}

\textbf{For Analysts:}
\begin{enumerate}[nosep]
\item Build segment profiles for your target population
\item Estimate complementarity $\gamma_{ij}$ from historical data
\item Design coherent portfolios covering all journey phases
\end{enumerate}

\end{tcolorbox}

% -----------------------------------------------------------------------------
% APPENDIX: QUICK REFERENCE
% -----------------------------------------------------------------------------
\newpage
\section*{Quick Reference Card}
\addcontentsline{toc}{section}{Quick Reference Card}

\begin{tcolorbox}[colback=lightgray, colframe=primaryblue, title=\textbf{Light Mode Checklist (F1--F6)}]
\small
\begin{tabular}{@{}rp{10cm}@{}}
$\square$ \textbf{F1} & Title: Clear, descriptive name \\
$\square$ \textbf{F2} & Type: One of 9 (Informative/Normative/Structural/...) \\
$\square$ \textbf{F3} & Context: Individual / Organization / Region / Nation \\
$\square$ \textbf{F4} & Target: IND / IDN / Collective utility \\
$\square$ \textbf{F5} & FEPSDE: Which dimension(s)? \\
$\square$ \textbf{F6} & Journey: Awareness / Trigger / Action / Maintenance / Stabilization \\
\end{tabular}
\end{tcolorbox}

\vspace{1em}

\begin{tcolorbox}[colback=lightgray, colframe=warningorange, title=\textbf{Axiom Violation Checklist}]
\small
\begin{tabular}{@{}rp{10cm}@{}}
$\square$ & Is every required field specified? (HHH-1) \\
$\square$ & Is context level explicit? (HHH-2) \\
$\square$ & Does intervention type match journey phase? (HHH-3) \\
$\square$ & Are there conflicting interventions ($\gamma < 0$)? (HHH-4) \\
$\square$ & Will this backfire on any segment? (HHH-5) \\
\end{tabular}
\end{tcolorbox}

\vspace{1em}

\begin{tcolorbox}[colback=lightgray, colframe=accentgreen, title=\textbf{Journey-Intervention Quick Match}]
\small
\begin{tabular}{@{}ll@{}}
\textbf{Awareness} & $\rightarrow$ Informative, Normative \\
\textbf{Trigger} & $\rightarrow$ Commitment, Temporal, Incentive \\
\textbf{Action} & $\rightarrow$ Structural, Feedback \\
\textbf{Maintenance} & $\rightarrow$ Feedback, Normative \\
\textbf{Stabilization} & $\rightarrow$ Identity-based, Structural \\
\end{tabular}
\end{tcolorbox}

% -----------------------------------------------------------------------------
% FOOTER
% -----------------------------------------------------------------------------
\vfill
\begin{center}
\textcolor{gray}{\small
Based on EBF Appendix HHH (METHOD-TOOLKIT)\\
For technical details, see the full appendix documentation.\\[0.5em]
\textcopyright{} 2026 FehrAdvice \& Partners AG
}
\end{center}

\end{document}
